% ============================================================
%  SECTION 2: CẤU HÌNH
% ============================================================
\chapter{Cấu hình}

\section{Chuẩn bị môi trường}

\subsection{Lý do chọn K3s thay Minikube}

Ban đầu đồ án được thực hiện trên Minikube. Sau quá trình triển khai, nhóm quyết định
chuyển sang \textbf{K3s v1.32.5} vì các lợi thế vượt trội:

\begin{table}[H]
	\centering
	\caption{So sánh Minikube và K3s}
	\begin{tabularx}{\textwidth}{L{4cm}L{4.5cm}L{4.5cm}}
		\toprule
		\textbf{Tiêu chí}         & \textbf{Minikube}                   & \textbf{K3s}                          \\
		\midrule
		Cơ chế chạy               & Docker container (Docker-in-Docker) & Systemd service (native process)      \\
		RAM overhead              & $\sim$2--3 GB                       & $\sim$512 MB                          \\
		Port 80/443               & Cần \texttt{minikube tunnel}        & Tự động (klipper-lb)                  \\
		Auto-start sau reboot     & Không                               & Có (\texttt{systemctl enable k3s})    \\
		inotify (Promtail, Kafka) & Phải set lại sau restart            & Persist trong \texttt{/etc/sysctl.d/} \\
		Worker node               & Docker container giả lập            & Máy thật (VM/máy khác)                \\
		\bottomrule
	\end{tabularx}
\end{table}

\subsection{Cài đặt K3s}

\begin{lstlisting}[style=bashstyle, caption=Cài K3s v1.32.5 (disable Traefik)]
# Cai K3s v1.32.5 stable LTS, tat Traefik (dung ingress-nginx thay the)
curl -sfL https://get.k3s.io | INSTALL_K3S_VERSION=v1.32.5+k3s1 sh -s - \
  --disable=traefik \
  --write-kubeconfig-mode=644

# Cau hinh kubeconfig
mkdir -p ~/.kube
sudo cp /etc/rancher/k3s/k3s.yaml ~/.kube/config
sudo chown $USER:$USER ~/.kube/config

# Kiem tra
kubectl get nodes
# NAME     STATUS   ROLES                  AGE   VERSION
# fedora   Ready    control-plane,master   1m    v1.32.5+k3s1
\end{lstlisting}

\begin{warnbox}[Lưu ý quan trọng về version]
	Không dùng K3s v1.35.x — có bug \texttt{cloud-controller-manager} exit ngay khi gặp 403 Forbidden
	lúc bootstrap RBAC (race condition), khiến K3s crash loop vĩnh viễn. Dùng v1.32.5 (LTS stable).
\end{warnbox}

\subsection{Cài Ingress NGINX và cấu hình inotify}

\begin{lstlisting}[style=bashstyle, caption=Cài ingress-nginx và tăng giới hạn inotify]
# Cai ingress-nginx
kubectl apply -f https://raw.githubusercontent.com/kubernetes/ingress-nginx/\
controller-v1.11.0/deploy/static/provider/cloud/deploy.yaml

# Tang gioi han inotify (can cho Kafka, Promtail)
sudo sysctl -w fs.inotify.max_user_watches=524288
sudo sysctl -w fs.inotify.max_user_instances=512

# Persist qua reboot (chi can set 1 lan voi K3s)
echo "fs.inotify.max_user_watches=524288" | sudo tee -a /etc/sysctl.d/99-k3s.conf
echo "fs.inotify.max_user_instances=512"  | sudo tee -a /etc/sysctl.d/99-k3s.conf
sudo sysctl --system
\end{lstlisting}

\subsection{Thêm Worker Node vào Cluster}

\begin{lstlisting}[style=bashstyle, caption=Join worker node vào K3s cluster]
# Tren control plane: lay token
sudo cat /var/lib/rancher/k3s/server/node-token

# Mo firewall tren control plane
sudo firewall-cmd --permanent --add-port=6443/tcp
sudo firewall-cmd --permanent --add-port=8472/udp
sudo firewall-cmd --reload

# Tren worker: cai K3s agent
curl -sfL https://get.k3s.io | INSTALL_K3S_VERSION=v1.32.5+k3s1 \
  K3S_URL="https://${CONTROL_PLANE_IP}:6443" \
  K3S_TOKEN="${TOKEN}" sh -

# Kiem tra tren control plane
kubectl get nodes
# NAME        STATUS   ROLES                  AGE
# fedora      Ready    control-plane,master   2d
# worker-01   Ready    <none>                 5m
# worker-02   Ready    <none>                 3m
\end{lstlisting}

% ============================================================
\section{Flow Setup — CI/CD Pipeline}

\subsection{Cấu hình GitHub Secrets}

Ba secrets bắt buộc cần cấu hình trong GitHub repository:

\begin{table}[H]
	\centering
	\caption{GitHub Secrets}
	\begin{tabularx}{\textwidth}{L{3.5cm}X}
		\toprule
		\textbf{Secret}       & \textbf{Mô tả}                                       \\
		\midrule
		\texttt{DOCKER\_USER} & Username Docker Hub — dùng để đặt tên image và login \\
		\texttt{DOCKER\_PASS} & Access Token Docker Hub — dùng để push image         \\
		\texttt{KUBE\_CONFIG} & Kubeconfig base64 — dùng để kubectl deploy lên K3s   \\
		\bottomrule
	\end{tabularx}
\end{table}

\begin{lstlisting}[style=bashstyle, caption=Lấy KUBE\_CONFIG cho K3s]
# Encode kubeconfig thanh base64
cat ~/.kube/config | base64 -w 0
# Copy output -> GitHub Settings -> Secrets -> KUBE_CONFIG
\end{lstlisting}

\subsection{Self-hosted Runner}

Vì K3s API server chạy tại \texttt{127.0.0.1:6443}, GitHub-hosted runner trên cloud
không thể kết nối được. Nhóm cài \textbf{self-hosted runner} trực tiếp trên máy chủ:

\begin{lstlisting}[style=bashstyle, caption=Cài và chạy self-hosted runner]
mkdir -p ~/actions-runner && cd ~/actions-runner
curl -o actions-runner-linux-x64-2.322.0.tar.gz -L \
  https://github.com/actions/runner/releases/download/v2.322.0/\
actions-runner-linux-x64-2.322.0.tar.gz
tar xzf ./actions-runner-linux-x64-2.322.0.tar.gz
./config.sh --url https://github.com/NPT-102/yas-2 --token <TOKEN>

# Chay background
nohup ./run.sh > runner.log 2>&1 &
\end{lstlisting}

\subsection{Workflow CI (Yêu cầu 3)}

File: \texttt{.github/workflows/npt-ci.yml}

\textbf{Luồng hoạt động:}
\begin{enumerate}
	\item Trigger: push lên bất kỳ branch nào (\texttt{branches: '**'})
	\item \textbf{Job detect-changes}: So sánh \texttt{git diff} để phát hiện service nào thay đổi
	\item \textbf{Job build-java}: Chạy \texttt{mvn clean install}, build Docker image, push lên Docker Hub
	      với tag = 7 ký tự đầu commit SHA
	\item \textbf{Job build-frontend}: Build Next.js qua multi-stage Dockerfile, push image
	\item Nếu push vào \texttt{main}: tag thêm \texttt{:latest}
\end{enumerate}

\begin{lstlisting}[style=yamlstyle, caption=CI workflow — detect changes và build (trích)]
name: CI Build and Push
on:
  push:
    branches: ['**']

jobs:
  detect-changes:
    runs-on: ubuntu-latest
    outputs:
      java_services: ${{ steps.filter.outputs.java_services }}
    steps:
      - uses: actions/checkout@v4
        with:
          fetch-depth: 2
      - name: Detect changed services
        id: filter
        run: |
          CHANGED=$(git diff --name-only HEAD~1 HEAD)
          for svc in $JAVA_SERVICES; do
            if echo "$CHANGED" | grep -q "^${svc}/"; then
              JAVA+=("\"${svc}\"")
            fi
          done
          echo "java_services=[$(IFS=,; echo "${JAVA[*]}")]" >> $GITHUB_OUTPUT

  build-java:
    needs: detect-changes
    if: needs.detect-changes.outputs.java_services != '[]'
    strategy:
      matrix:
        service: ${{ fromJson(needs.detect-changes.outputs.java_services) }}
    steps:
      - name: Build and Push Docker Image
        run: |
          COMMIT_ID=${GITHUB_SHA::7}
          IMAGE=${{ secrets.DOCKER_USER }}/${{ matrix.service }}
          TAGS="-t ${IMAGE}:${COMMIT_ID}"
          if [ "${{ github.ref_name }}" = "main" ]; then
            TAGS="$TAGS -t ${IMAGE}:latest"
          fi
          docker buildx build --push $TAGS ./${{ matrix.service }}
\end{lstlisting}

\subsection{Workflow Developer Build (Yêu cầu 4)}

File: \texttt{.github/workflows/npt-developer\_build.yml}

Developer kích hoạt thủ công qua \texttt{workflow\_dispatch}, nhập:
\begin{itemize}
	\item \texttt{service\_name}: service muốn test (dropdown list)
	\item \texttt{branch\_name}: tên branch đang phát triển
\end{itemize}

\textbf{Luồng hoạt động:}
\begin{enumerate}
	\item Lấy commit SHA cuối của branch được chọn
	\item Checkout code của branch đó \textrightarrow{} build + push image với tag = SHA
	\item Deploy tất cả service vào namespace \texttt{developer}:
	      \begin{itemize}
		      \item Service được chọn: dùng image mới build, expose qua NodePort
		      \item Các service còn lại: dùng image \texttt{:latest}
	      \end{itemize}
	\item In ra GitHub Step Summary: NodePort, Worker IP, hướng dẫn \texttt{/etc/hosts}
\end{enumerate}

\begin{lstlisting}[style=yamlstyle, caption=Developer build — deploy service với custom image (trích)]
- name: Deploy all services
  run: |
    for svc in $SERVICES; do
      if [ "${svc}" = "${CUSTOM_SERVICE}" ]; then
        # Deploy custom branch image
        helm upgrade --install ${svc} ./${svc} \
          --namespace developer --create-namespace \
          --set backend.image.repository=${DOCKER_USER}/${svc} \
          --set backend.image.tag=${CUSTOM_TAG} \
          --set backend.service.type=NodePort \
          --set backend.ingress.host="api.developer.yas.local.com"
      else
        # Deploy latest image
        helm upgrade --install ${svc} ./${svc} \
          --namespace developer --create-namespace \
          --set backend.image.tag=latest
      fi
    done
\end{lstlisting}

\subsection{Workflow Cleanup (Yêu cầu 5)}

File: \texttt{.github/workflows/npt-cleanup.yml}

Developer chạy workflow, nhập \texttt{yes} để xác nhận \textrightarrow{} xóa toàn bộ namespace \texttt{developer}.

\begin{lstlisting}[style=bashstyle, caption=Cleanup deployment]
kubectl delete namespace developer --ignore-not-found=true
\end{lstlisting}

\subsection{Workflow Dev Deploy (Yêu cầu 6a)}

File: \texttt{.github/workflows/npt-dev-deploy.yml}

Trigger: \textbf{tự động} khi push vào \texttt{main}. Deploy tất cả service vào namespace \texttt{dev}.
ArgoCD ApplicationSet đồng bộ Helm chart từ Git \textrightarrow{} tự động sync khi \texttt{main} thay đổi.

\subsection{Workflow Staging Deploy (Yêu cầu 6b)}

File: \texttt{.github/workflows/npt-staging-deploy.yml}

Trigger: tạo tag \texttt{v*} trên branch \texttt{main}.
\begin{lstlisting}[style=bashstyle, caption=Tạo staging release]
git checkout main && git pull
git tag v1.0.0
git push origin v1.0.0
# Pipeline tu dong build tat ca service voi tag v1.0.0
# Push images: <DOCKER_USER>/<service>:v1.0.0
# Deploy vao namespace staging
\end{lstlisting}

% ============================================================
\section{Cấu hình ArgoCD (Nâng cao)}

\subsection{Cài đặt ArgoCD}

\begin{lstlisting}[style=bashstyle, caption=Cài ArgoCD với server-side apply]
kubectl create namespace argocd
kubectl apply -n argocd \
  -f https://raw.githubusercontent.com/argoproj/argo-cd/stable/manifests/install.yaml \
  --server-side  # Can thiet vi CRD ArgoCD lon hon 256KB

# Truy cap UI qua port-forward
kubectl port-forward svc/argocd-server -n argocd 9090:443 &
# https://localhost:9090  (admin / <mat-khau-lay-tu-secret>)
\end{lstlisting}

\subsection{Cấu hình ApplicationSet}

ArgoCD ApplicationSet tự động tạo Application cho từng service trong mỗi environment.

\begin{lstlisting}[style=yamlstyle, caption=ArgoCD ApplicationSet cho môi trường dev]
apiVersion: argoproj.io/v1alpha1
kind: ApplicationSet
metadata:
  name: yas-dev-appset
  namespace: argocd
spec:
  generators:
    - list:
        elements:
          - service: cart
          - service: order
          - service: product
          # ... them 17 services
  template:
    metadata:
      name: 'yas-dev-\{\{service\}\}'
    spec:
      project: yas
      source:
        repoURL: https://github.com/NPT-102/yas-2
        targetRevision: main
        path: 'k8s/charts/\{\{service\}\}'
      destination:
        server: https://kubernetes.default.svc
        namespace: dev
      syncPolicy:
        automated:
          prune: true
          selfHeal: true
        syncOptions:
          - CreateNamespace=true
\end{lstlisting}

% ============================================================
\section{Cấu hình Istio Service Mesh (Nâng cao)}

\subsection{Cài đặt Istio}

\begin{lstlisting}[style=bashstyle, caption=Cài Istio với profile tối ưu]
# Cai Istio voi resource limits toi uu (istio-overlay.yaml)
./istio-1.24.3/bin/istioctl install -f istio/istio-overlay.yaml -y

# Bat sidecar injection
kubectl label namespace yas istio-injection=enabled
kubectl label namespace ingress-nginx istio-injection=enabled  # QUAN TRONG!

# Restart de nhan sidecar (2/2 containers)
kubectl rollout restart deployment -n yas
kubectl rollout restart deployment ingress-nginx-controller -n ingress-nginx

# Cai addons
kubectl apply -f https://raw.githubusercontent.com/istio/istio/release-1.24/samples/addons/prometheus.yaml
kubectl apply -f https://raw.githubusercontent.com/istio/istio/release-1.24/samples/addons/grafana.yaml
kubectl apply -f https://raw.githubusercontent.com/istio/istio/release-1.24/samples/addons/kiali.yaml
\end{lstlisting}

\subsection{Giải thích code cấu hình Istio}

\subsubsection{PeerAuthentication — bắt buộc mTLS}

\begin{lstlisting}[style=yamlstyle, caption=istio/peer-authentication.yaml]
apiVersion: security.istio.io/v1beta1
kind: PeerAuthentication
metadata:
  name: default-mtls
  namespace: yas
spec:
  mtls:
    mode: STRICT  # Bat buoc mTLS cho toan bo traffic vao namespace yas
                  # ingress-nginx co sidecar -> tu dung mTLS qua DestinationRule
\end{lstlisting}

\textbf{Giải thích:} \texttt{STRICT} mode yêu cầu mọi kết nối vào namespace \texttt{yas}
phải dùng mTLS. Điều này có nghĩa là chỉ các client \textit{trong mesh} (có Envoy sidecar)
mới giao tiếp được. Đây là lý do ingress-nginx cũng phải được inject sidecar.

\subsubsection{DestinationRule — cấu hình TLS mode}

\begin{lstlisting}[style=yamlstyle, caption=istio/destination-rule.yaml (trích — 1 trong 21 rules)]
apiVersion: networking.istio.io/v1beta1
kind: DestinationRule
metadata:
  name: mtls-cart
  namespace: yas
spec:
  host: cart.yas.svc.cluster.local
  exportTo:
    - "."           # Namespace yas
    - "ingress-nginx"  # Quan trong: ingress-nginx phai thay duoc rule nay
  trafficPolicy:
    tls:
      mode: ISTIO_MUTUAL  # Envoy tu quan ly chung chi, khong can config thu cong
\end{lstlisting}

\textbf{Giải thích:} \texttt{exportTo: ingress-nginx} cho phép sidecar trong namespace
\texttt{ingress-nginx} nhìn thấy DestinationRule và tự dùng mTLS. Nếu thiếu field này,
ingress-nginx proxy qua Pod IP (PassthroughCluster) thay vì service identity \textrightarrow{} mTLS STRICT fail.

\subsubsection{AuthorizationPolicy — kiểm soát service-to-service}

\begin{lstlisting}[style=yamlstyle, caption=istio/authorization-policy.yaml (trích)]
# Policy mac dinh: deny all
apiVersion: security.istio.io/v1beta1
kind: AuthorizationPolicy
metadata:
  name: deny-all
  namespace: yas
spec: {}   # Khong co rules = deny tat ca

---
# Chi cho phep order truy cap tax
apiVersion: security.istio.io/v1beta1
kind: AuthorizationPolicy
metadata:
  name: allow-order-to-tax
  namespace: yas
spec:
  selector:
    matchLabels:
      app.kubernetes.io/name: tax   # Pod nao ap dung policy nay
  action: ALLOW
  rules:
    - from:
        - source:
            principals:
              - "cluster.local/ns/yas/sa/order"  # Chi order duoc phep
    - from:
        - source:
            principals:
              - "cluster.local/ns/ingress-nginx/sa/ingress-nginx"
\end{lstlisting}

\subsubsection{VirtualService — Retry Policy}

\begin{lstlisting}[style=yamlstyle, caption=istio/virtual-service-retry.yaml (trích)]
apiVersion: networking.istio.io/v1beta1
kind: VirtualService
metadata:
  name: tax-retry
  namespace: yas
spec:
  hosts:
    - tax.yas.svc.cluster.local
  http:
    - retries:
        attempts: 3          # Thu lai 3 lan neu loi
        perTryTimeout: 10s
        retryOn: 5xx,reset,connect-failure,retriable-4xx
      route:
        - destination:
            host: tax.yas.svc.cluster.local
\end{lstlisting}

\textbf{Giải thích:} Nếu service \texttt{tax} trả lỗi 5xx, Istio tự động retry tối đa 3 lần
mà không cần code application xử lý retry. Điều này tăng độ bền của hệ thống.

\subsection{Ingress NGINX kết hợp Istio mTLS STRICT}

Thách thức lớn nhất: \texttt{ingress-nginx} ban đầu nằm ngoài mesh \textrightarrow{} gửi plain HTTP
vào namespace \texttt{yas} \textrightarrow{} sidecar từ chối (STRICT mode) \textrightarrow{} 502 Bad Gateway.

\textbf{Giải pháp 3 bước:}
\begin{enumerate}
	\item Inject sidecar vào \texttt{ingress-nginx} (label namespace + restart pod)
	\item Export DestinationRule sang namespace \texttt{ingress-nginx}
	\item Thêm annotation vào Ingress để proxy theo service identity (không theo Pod IP):
\end{enumerate}

\begin{lstlisting}[style=yamlstyle, caption=Annotation Ingress để ingress-nginx dùng service identity]
# k8s/charts/backend/templates/ingress.yaml
annotations:
  nginx.ingress.kubernetes.io/service-upstream: "true"
  nginx.ingress.kubernetes.io/upstream-vhost: "{{ .Chart.Name }}.yas.svc.cluster.local"
\end{lstlisting}

Khi bật \texttt{service-upstream}, Envoy sidecar của ingress-nginx route traffic qua
cluster name \texttt{outbound|80||cart.yas.svc.cluster.local} thay vì
\texttt{PassthroughCluster} \textrightarrow{} DestinationRule được áp dụng \textrightarrow{} mTLS hoạt động.

% ============================================================
\section{Cấu hình Infrastructure}

\subsection{Setup cluster script}

Script \texttt{k8s/deploy/setup-cluster.sh} tự động cài toàn bộ infrastructure:

\begin{lstlisting}[style=bashstyle, caption=Chạy setup cluster]
cd k8s/deploy
./setup-cluster.sh    # ~15 phut
./setup-keycloak.sh   # ~5 phut
./setup-redis.sh      # ~1 phut
./deploy-yas-configuration.sh  # ConfigMaps + Secrets
./deploy-yas-applications.sh   # Deploy core services + mesh
\end{lstlisting}

\subsection{Cấu hình CoreDNS cho hostname nội bộ}

Pods trong cluster không đọc được \texttt{/etc/hosts} của host \textrightarrow{} cần thêm hostname
vào CoreDNS:

\begin{lstlisting}[style=bashstyle, caption=Thêm hostname vào CoreDNS]
kubectl edit configmap coredns -n kube-system
# Them vao block "hosts /etc/coredns/NodeHosts {"
# 172.16.0.240 identity.yas.local.com
# 172.16.0.240 api.yas.local.com

kubectl rollout restart deployment/coredns -n kube-system
\end{lstlisting}
